\documentclass[11pt,a4paper]{article}

% --- Mathematics (must precede unicode-math) ---
\usepackage{amsmath}
\usepackage{mathtools}

% --- Fonts (lualatex) ---
\usepackage{fontspec}
\usepackage{unicode-math}
\setmathfont{Latin Modern Math}

% --- Microtypography (after all font setup) ---
\usepackage{microtype}

% --- Colors & page layout ---
\usepackage[dvipsnames]{xcolor}
\usepackage[margin=25mm, top=30mm, bottom=30mm]{geometry}

% --- Headers & footers ---
\usepackage{fancyhdr}
\pagestyle{fancy}
\fancyhf{}
\fancyhead[L]{\small\itshape Near-Field Green's Tensor Verification}
\fancyhead[R]{\small\itshape Wu \& Ben-Menahem (1985)}
\fancyfoot[C]{\thepage}
\renewcommand{\headrulewidth}{0.4pt}
\renewcommand{\footrulewidth}{0pt}

% --- Section numbering ---
\setcounter{secnumdepth}{3}

% --- Tables ---
\usepackage{booktabs}
\usepackage{array}

% --- Captions ---
\usepackage[font=small, labelfont=bf, labelsep=period]{caption}

% --- Spacing ---
\usepackage{setspace}
\setstretch{1.15}
\setlength{\parskip}{0.4em}

% --- Citations ---
\usepackage[round,authoryear]{natbib}
\bibliographystyle{plainnat}

% --- Hyperlinks (last before macros) ---
\usepackage[colorlinks=true, linkcolor=blue!60!black,
            citecolor=blue!60!black, urlcolor=blue!60!black]{hyperref}

% --- Macros ---
\newcommand{\bG}{\mathbf{G}}
\newcommand{\bI}{\mathbf{I}}
\newcommand{\bx}{\mathbf{x}}
\newcommand{\beR}{\hat{\mathbf{e}}_R}
\newcommand{\dd}{\mathrm{d}}
\newcommand{\order}[1]{\mathcal{O}\!\left(#1\right)}

\begin{document}

% --- Title block ---
\thispagestyle{plain}
\begin{center}
\vspace*{1em}
{\LARGE\bfseries Derivation and Verification of the\\[0.3em]
Elastodynamic Near-Field Green's Tensor}\\[1.8em]
{\large T.\,M.\ Nestor}\\[0.5em]
{\itshape Research School of Earth Sciences, Australian National University}\\[2em]
\end{center}

\section{Introduction}

The displacement field due to a point force in a homogeneous isotropic
elastic medium is governed by the Stokes tensor, first derived by
\citet{Stokes1849}. In the frequency domain, this tensor exhibits a
rich structure with contributions at different powers of the
source--receiver distance~$R$: a far-field term ($1/R$), and
near-field terms ($1/R^2$ and~$1/R^3$).

\citet{WuBenMenahem1985} showed that in the near-source region
($k_\beta R \ll 1$), the $1/R^3$ and $1/R^2$ contributions from
P-waves and S-waves cancel due to their mutual coupling through the
equation of motion. The surviving residue at $\order{1/R}$ is a
distinctive quadrupolar term proportional to
$(\bI - 3\,\beR\beR^{\!\top})$, controlled by the velocity contrast
factor $(1 - \beta^2/\alpha^2)$.

This document presents a self-contained derivation of this result
starting from the Kupradze representation
\citep{Kupradze1963,Kausel2006} and provides both symbolic
(Mathematica) and numerical (Python) verification.

\section{The Kupradze Representation}

The exact frequency-domain elastodynamic Green's tensor satisfies
\citep{AkiRichards2002,BenMenahemSingh1981}:
\begin{equation}
  \mu\,\nabla^2 G_{ij}
  + (\lambda + \mu)\,\partial_i\partial_k G_{kj}
  + \rho\omega^2 G_{ij}
  = -\delta_{ij}\,\delta(\bx).
  \label{eq:pde}
\end{equation}
Its solution admits the Kupradze representation
\citep{Kupradze1963}:
\begin{equation}
  \boxed{
  G_{ij}(\bx,\omega)
  = \frac{1}{4\pi\mu}
    \biggl[
      \delta_{ij}\,\Phi_\beta
      + \frac{1}{k_\beta^2}\,
        \partial_i\partial_j\!\left(\Phi_\beta - \Phi_\alpha\right)
    \biggr],
  }
  \label{eq:kupradze}
\end{equation}
where
$\Phi_k \equiv \Phi(k,R) = e^{-ikR}/R$
is the outgoing scalar Helmholtz Green's function satisfying
$(\nabla^2 + k^2)\Phi_k = -4\pi\,\delta(\bx)$, and
\begin{equation}
  k_\alpha = \frac{\omega}{\alpha}, \quad
  k_\beta = \frac{\omega}{\beta}, \quad
  \alpha^2 = \frac{\lambda + 2\mu}{\rho}, \quad
  \beta^2 = \frac{\mu}{\rho}.
  \label{eq:wavenumbers}
\end{equation}

\subsection{Verification of the Kupradze representation}

We verify \eqref{eq:kupradze} by direct substitution into~\eqref{eq:pde}.
Using $\nabla^2\Phi_k = -k^2\Phi_k - 4\pi\,\delta(\bx)$, each term evaluates as:
\begin{align}
  \mu\,\nabla^2 G_{ij}
  &= \frac{1}{4\pi}\!\left[
       -k_\beta^2\,\delta_{ij}\,\Phi_\beta - 4\pi\,\delta_{ij}\,\delta(\bx)
       - \frac{1}{k_\beta^2}\,\partial_i\partial_j
         \bigl(k_\beta^2\Phi_\beta - k_\alpha^2\Phi_\alpha\bigr)
     \right], \label{eq:term1} \\[4pt]
  (\lambda\!+\!\mu)\,\partial_i\partial_k G_{kj}
  &= \frac{\lambda\!+\!\mu}{4\pi\mu}\,
     \frac{k_\alpha^2}{k_\beta^2}\,\partial_i\partial_j\Phi_\alpha,
     \label{eq:term2} \\[4pt]
  \rho\omega^2 G_{ij}
  &= \frac{k_\beta^2}{4\pi}\!\left[
       \delta_{ij}\,\Phi_\beta
       + \frac{1}{k_\beta^2}\,\partial_i\partial_j
         \bigl(\Phi_\beta - \Phi_\alpha\bigr)
     \right]. \label{eq:term3}
\end{align}
In the sum~\eqref{eq:term1}+\eqref{eq:term2}+\eqref{eq:term3}, the
$\delta_{ij}\Phi_\beta$ terms cancel, the $\partial_i\partial_j\Phi_\beta$
terms cancel, and the $\partial_i\partial_j\Phi_\alpha$ coefficient
vanishes via
$k_\alpha^2/k_\beta^2 = \mu/(\lambda+2\mu)$, yielding
$-\delta_{ij}\,\delta(\bx)$ as required.~$\square$

\section{Taylor Expansion in the Near-Source Limit}

\subsection{Second derivatives of the scalar Green's function}

The second spatial derivative of $\Phi_k = e^{-ikR}/R$ decomposes on
the direction dyadic as:
\begin{equation}
  \partial_i\partial_j\Phi_k
  = e^{-ikR}\!\left[
      \left(\frac{-k^2}{R} + \frac{3ik}{R^2} + \frac{3}{R^3}\right)
      \gamma_i\gamma_j
      + \left(\frac{-ik}{R^2} - \frac{1}{R^3}\right)\delta_{ij}
    \right],
  \label{eq:d2phi}
\end{equation}
where $\gamma_i = x_i/R$ is the direction cosine.

\subsection{Expansion of the second-derivative difference}

Expanding the exponentials $e^{-ikR} = 1 - ikR - k^2R^2/2 + \cdots$
and collecting terms order-by-order in~$R$:

\paragraph{The \texorpdfstring{$1/R^3$}{1/R3} terms.}
\begin{equation}
  \frac{1}{R^3}\bigl[(3\gamma_i\gamma_j - \delta_{ij})(1-1)\bigr] = 0.
\end{equation}
The unit leading coefficients cancel identically between $\Phi_\beta$
and~$\Phi_\alpha$, since both share the same $1/R^3$ structure.

\paragraph{The \texorpdfstring{$1/R^2$}{1/R2} terms.}
From the $(-ikR)$ part of the exponential hitting the $1/R^3$ piece,
and the leading part of the $ik/R^2$ piece:
\begin{equation}
  \frac{1}{R^2}\bigl[
    (3\gamma_i\gamma_j - \delta_{ij})
    \bigl(-i(k_\beta - k_\alpha) + i(k_\beta - k_\alpha)\bigr)
  \bigr] = 0.
\end{equation}
Again, exact cancellation between cross-terms.

\paragraph{The surviving \texorpdfstring{$1/R$}{1/R} term.}
The first non-vanishing contribution comes at $\order{1/R}$.
Collecting all terms:
\begin{equation}
  \partial_i\partial_j(\Phi_\beta - \Phi_\alpha)
  = -\frac{k_\beta^2 - k_\alpha^2}{2R}\,
    (\delta_{ij} - \gamma_i\gamma_j)
  + \order{R^0}.
  \label{eq:d2diff_asymp}
\end{equation}

\subsection{Complete near-source asymptotic}

Substituting~\eqref{eq:d2diff_asymp} and $\Phi_\beta \approx 1/R$
into the Kupradze representation~\eqref{eq:kupradze}:
\begin{align}
  G_{ij}
  &\approx \frac{1}{4\pi\mu R}\,\delta_{ij}
    - \frac{k_\beta^2 - k_\alpha^2}{8\pi\mu\,k_\beta^2\,R}\,
      (\delta_{ij} - \gamma_i\gamma_j) \notag\\[4pt]
  &= \frac{1}{4\pi\mu R}
     \biggl[\delta_{ij}
       - \frac{1}{2}\!\left(1 - \frac{\beta^2}{\alpha^2}\right)
         (\delta_{ij} - \gamma_i\gamma_j)
     \biggr],
  \label{eq:full_asymp}
\end{align}
using
\begin{equation}
  \frac{k_\beta^2 - k_\alpha^2}{k_\beta^2}
  = 1 - \frac{k_\alpha^2}{k_\beta^2}
  = 1 - \frac{\mu}{\lambda + 2\mu}
  = 1 - \frac{\beta^2}{\alpha^2}
  = \frac{\lambda + \mu}{\lambda + 2\mu}.
  \label{eq:contrast}
\end{equation}
Expanding and simplifying~\eqref{eq:full_asymp} yields the compact form:

\begin{equation}
  \boxed{
  \bG \;\underset{R\to 0}{\approx}\;
  \frac{1}{8\pi\mu R}
  \biggl[
    \left(1 + \frac{\beta^2}{\alpha^2}\right)\bI
    + \left(1 - \frac{\beta^2}{\alpha^2}\right)\beR\beR^{\!\top}
  \biggr].
  }
  \label{eq:eq6}
\end{equation}
This is Eq.~(6) of \citet{WuBenMenahem1985}.

\section{Decomposition into Near-Field and Far-Field Terms}

Following \citet{WuBenMenahem1985}, the complete
asymptotic~\eqref{eq:eq6} can be decomposed as:
\begin{equation}
  \bG \approx \underbrace{
    \frac{-1}{8\pi\mu R}
    \left(1 - \frac{\beta^2}{\alpha^2}\right)
    (\bI - 3\,\beR\beR^{\!\top})
  }_{\displaystyle (\bG)_\text{near} \;\text{[Eq.~(5)]}}
  \;+\;
  \underbrace{
    \frac{1}{4\pi\mu R}
    \left[
      (\bI - \beR\beR^{\!\top})
      + \frac{\beta^2}{\alpha^2}\,\beR\beR^{\!\top}
    \right]
  }_{\displaystyle (\bG)_\text{far} \;\text{[leading order]}}.
  \label{eq:decomp}
\end{equation}

The first term is the ``near-field term'' of \citet{WuBenMenahem1985}
--- the $\order{1/R}$ residue from the mutual cancellation of the
P-wave and S-wave static-like ($1/R^2$, $1/R^3$) contributions. It has the
distinctive \emph{quadrupolar} radiation pattern
$(\bI - 3\,\beR\beR^{\!\top})$, which is traceless and vanishes
when~$\alpha = \beta$.

The second term represents the far-field ($1/R$ propagating) P and S
contributions evaluated at their leading order in the near-source
limit. It carries the standard transverse/longitudinal projector structure.

\begin{quote}
\textbf{Key point:} The ``near-field term'' alone, Eq.~(5) of
\citet{WuBenMenahem1985},
\begin{equation}
  (\bG)_\text{near}
  = \frac{-1}{8\pi\mu R}
    \left(1 - \frac{\beta^2}{\alpha^2}\right)
    (\bI - 3\,\beR\beR^{\!\top}),
  \label{eq:eq5}
\end{equation}
is \emph{not} the complete near-source Green's tensor. The full
$R \to 0$ limit also includes the leading-order far-field
contribution. As confirmed numerically (Table~\ref{tab:convergence}),
Eq.~\eqref{eq:eq5} alone produces ${\sim}105\%$ relative error, while
the full asymptotic~\eqref{eq:eq6} converges as~$\order{k_\beta R}$.
\end{quote}

\section{Cancellation Mechanism}

The cancellation of the $1/R^3$ and $1/R^2$ terms deserves further
comment. From the Kupradze form~\eqref{eq:kupradze}, the near-source
singularity of the Green's tensor is governed by the difference
$\Phi_\beta - \Phi_\alpha$. Both P and S scalar Green's functions share
the \emph{same} static singularity
$\Phi_k \to 1/R$ as $kR \to 0$, so their difference
\begin{equation}
  \Phi_\beta - \Phi_\alpha
  = \frac{e^{-ik_\beta R} - e^{-ik_\alpha R}}{R}
  \;\xrightarrow{kR \to 0}\;
  \frac{-i(k_\beta - k_\alpha)R}{R} + \order{R}
  = -i(k_\beta - k_\alpha) + \order{R}
  \label{eq:phi_diff}
\end{equation}
starts at $\order{R^0}$, \emph{not} $\order{1/R}$.  Consequently, the
second derivative $\partial_i\partial_j(\Phi_\beta - \Phi_\alpha)$
starts at $\order{1/R}$ rather than the naïve $\order{1/R^3}$.  This
is the mathematical origin of the cancellation first noted by
\citet{Pendse1948} and analysed in detail by \citet{MorseFeshbach1953}
and \citet{WuBenMenahem1985}.

Physically, in the static limit both P and S displacements must
combine to satisfy the equilibrium equation $\nabla \cdot \symbf{\sigma}
= 0$. The near-source singularity is fixed by the \emph{total}
elastic response, not by P or S individually, forcing the static-like
parts to be in exact balance \citep{BenMenahemSingh1981,Harkrider1976}.

\section{Numerical Verification}

Table~\ref{tab:convergence} compares the exact Stokes tensor
(Kupradze representation, evaluated analytically) against the
near-field term~\eqref{eq:eq5} and the full
asymptotic~\eqref{eq:eq6}, for a granite-like medium with
$\alpha = 5292$\,m/s, $\beta = 3162$\,m/s, $\rho = 2700$\,kg/m$^3$
at 1\,Hz, along the observation direction
$\hat{\bx} = (1,1,1)/\sqrt{3}$.

\begin{table}[h!]
\centering
\caption{Relative error $\|\bG_\text{exact} - \bG_\text{approx}\|
  / \|\bG_\text{exact}\|$ for the two approximations.
  The Eq.~(6) error tracks the imaginary-to-real ratio,
  confirming $\order{k_\beta R}$ convergence.}
\label{tab:convergence}
\begin{tabular}{@{}rrccc@{}}
\toprule
$R/\lambda_S$ & $k_\beta R$ &
\multicolumn{1}{c}{Error, Eq.~(5)} &
\multicolumn{1}{c}{Error, Eq.~(6)} &
\multicolumn{1}{c}{$|\mathrm{Im}\,\bG|/|\mathrm{Re}\,\bG|$} \\
\midrule
$10^{-4}$ & 0.0006 & 105.2\% & 0.06\% & 0.06\% \\
$10^{-3}$ & 0.006  & 105.2\% & 0.6\%  & 0.6\%  \\
$5{\times}10^{-3}$ & 0.031 & 105.2\% & 2.9\% & 2.9\% \\
$10^{-2}$ & 0.063  & 105.2\% & 5.8\%  & 5.8\%  \\
$5{\times}10^{-2}$ & 0.31 & 104.9\% & 28.9\% & 29.8\% \\
$10^{-1}$ & 0.63   & 104.0\% & 57.0\% & 64.8\% \\
\bottomrule
\end{tabular}
\end{table}

\noindent
Additionally, the following symbolic verifications were performed
in Mathematica (see accompanying notebook
\texttt{NearField\_Verification.nb}):
\begin{enumerate}
  \item Taylor expansion of the exact Kupradze tensor to
        $\order{1/R}$ reproduces Eq.~\eqref{eq:eq6} identically.
  \item The $1/R^3$ and $1/R^2$ terms vanish (P--S cancellation).
  \item Eq.~\eqref{eq:eq6} decomposes as
        $\text{Eq.~\eqref{eq:eq5}} + (\bG)_\text{far}$.
  \item The identity
        $1 - \beta^2/\alpha^2 = (\lambda+\mu)/(\lambda+2\mu)$ holds.
\end{enumerate}
All four tests pass.

A numerical (finite-difference) implementation of the Kupradze
representation confirms the analytic formulae to $\order{10^{-5}}$
relative accuracy, providing an independent cross-check (see
\texttt{verify\_nearfield.py}).

\section{Connection to the T-Matrix Formulation}

In the T-matrix approach to elastic wave scattering
\citep{Nestor1996}, the Green's tensor enters the Lippmann--Schwinger
integral equation. For scatterers whose dimensions are small compared
to the wavelength ($ka \ll 1$, the Rayleigh regime), the
near-field form of the Green's tensor governs the
interaction between closely spaced inclusions. The
velocity contrast factor $(1 - \beta^2/\alpha^2)$ controls the
strength of the near-field coupling, vanishing in the acoustic limit.

\bigskip

\bibliography{nearfield_verification}

\end{document}
